% GENERATED FILE. Edit canonical lessons, then run npm run generate:books.
% canonical-source-hash: fnv1a64:55c871611b787d03
% canonical-lessons: AR-C292-itamada, AR-C292-yatamidu-ala, AR-C292-law, AR-R292-would-and-depends, AR-W292-sitta-sutur

\chapter{Would, And Six Lines About A Trip}
\label{ch:ar-would-and-six-lines}
\hlchaptermodality{292}

\begin{chapteropening}
\textbf{By the end of this chapter:} \emph{I can say what I would do if things were different, say that a plan depends on something, and write six connected lines about a trip.}

Complete AR-R292-would-and-depends: an unreal condition with law, a wish with layta, and a plan that depends on something, with ʿalā ḥasab and the condition word.
\end{chapteropening}

\section[\texorpdfstring{\ar{اعتمد} (iʿtamada)}{اعتمد (iʿtamada)}]{\ar{اعتمد} (iʿtamada) — to depend on}

\label{lesson:AR-C292-itamada}



\begin{quote}
Say \textquotedblleft{}if it rains, we'll watch a film\textquotedblright{}. (\emph{Idhā amṭarat, fa-sa-nushāhidu fīlman}.)

Sometimes you cannot give the result yet, because it hangs on something else. This is the verb for hanging on something.
\end{quote}

\subsection*{You'll want to know: \ar{اعتمد}}
\textbf{\ar{اعتمد}} — \emph{iʿtamada} — \textquotedblleft{}he depended, he relied\textquotedblright{}, with \textbf{\ar{على}}, \emph{ʿalā}, \textquotedblleft{}on\textquotedblright{}.

\begin{quote}
\textbf{\ar{اعتمدت} \ar{على} \ar{صديقي}.} — \emph{iʿtamadtu ʿalā ṣadīqī} — I relied on my friend.
\end{quote}

The present is \emph{yaʿtamidu}, he depends; \emph{aʿtamidu ʿalayka}, I'm counting on you. It has the same shape as \emph{iʿtaqada}, to believe, and its present \emph{aʿtaqidu}: one letter apart.

\begin{cousinweb}
\textbf{\ar{اعتمد}} (\emph{iʿtamada}) is built on the native Arabic root \emph{ʿ-m-d}, \textquotedblleft{}to prop up, to support\textquotedblright{} (Lane). Lane gives \emph{iʿtamadtu ʿalā ash-shayʾ} as leaning or resting on a thing, and then relying on a person. To depend on something, in Arabic, is to lean your weight on it.
\end{cousinweb}

\subsection*{Your turn}
Say these aloud:

\begin{itemize}
\raggedright
  \item \emph{iʿtamada}, then \emph{yaʿtamidu}
  \item \textquotedblleft{}I relied on my friend\textquotedblright{} — \emph{iʿtamadtu ʿalā ṣadīqī}
  \item \textquotedblleft{}I'm counting on you\textquotedblright{} — \emph{aʿtamidu ʿalayka}
  \item \textquotedblleft{}We don't depend on money\textquotedblright{} — \emph{lā naʿtamidu ʿalā al-māl}
\end{itemize}

\begin{tcolorbox}[breakable,colback=teal!4,colframe=teal!35!black,title={Before you move on}]
What does \ar{اعتمد} mean? (\textquotedblleft{}To depend on, to rely on\textquotedblright{}.) Which little word goes with it? (\emph{\textbf{ʿAlā}}, on.) Which verb has the same shape? (\emph{\textbf{Iʿtaqada}}, to believe.)
\end{tcolorbox}
\section[\texorpdfstring{\ar{يعتمد} \ar{على} (yaʿtamidu ʿalā)}{يعتمد على (yaʿtamidu ʿalā)}]{\ar{يعتمد} \ar{على} (yaʿtamidu ʿalā) — it depends on}

\label{lesson:AR-C292-yatamidu-ala}



\begin{quote}
Say \textquotedblleft{}I relied on my friend\textquotedblright{}. (\emph{Iʿtamadtu ʿalā ṣadīqī}.) Then \textquotedblleft{}if it rains, we won't go out\textquotedblright{}. (\emph{Idhā amṭarat, fa-lan nakhruja}.)

Someone asks about your plan, and you do not know yet. That is an answer too.
\end{quote}

\begin{grammarlens}[title={Grammar Lens: it depends on}]
\begin{quote}
— \textbf{\ar{هل} \ar{ستسافر} \ar{غدا؟}} — \emph{hal sa-tusāfiru ghadan?} — Are you travelling tomorrow? — \textbf{\ar{يعتمد} \ar{على} \ar{الطقس}.} — \emph{yaʿtamidu ʿalā aṭ-ṭaqs} — It depends on the weather.
\end{quote}

\textquotedblleft{}It\textquotedblright{} is not said: the \emph{ya-} of \emph{yaʿtamidu} carries it. Put the thing the plan hangs on after \emph{ʿalā}: \emph{yaʿtamidu ʿalā al-waqt}, it depends on the time; \emph{yaʿtamidu ʿalayka}, it's up to you.

It answers the same need as a condition with \emph{idhā}, but leaves the condition open: \emph{idhā kāna aṭ-ṭaqs jamīlan, fa-sa-usāfiru} says which way it goes; \emph{yaʿtamidu ʿalā aṭ-ṭaqs} only says what decides.
\end{grammarlens}

\begin{grammarlens}[title={Grammar Lens: ḥasaba, according to}]
\emph{Ḥasaba}, from chapter two hundred and sixty-eight, puts the deciding thing in front of a plain statement:

\begin{quote}
\textbf{\ar{حسب} \ar{الجدول،} \ar{ينطلق} \ar{القطار} \ar{بعد} \ar{عشر} \ar{دقائق}.} — \emph{ḥasaba al-jadwal, yanṭaliqu al-qiṭār baʿda ʿashr daqāʾiq} — According to the timetable, the train leaves in ten minutes.
\end{quote}

\emph{Ḥasaba aṭ-ṭaqs} is \textquotedblleft{}depending on the weather\textquotedblright{}: \emph{sa-nusāfiru bi-l-qiṭār aw bi-l-ḥāfila, ḥasaba aṭ-ṭaqs}.
\end{grammarlens}

\subsection*{Your turn}
Say these aloud:

\begin{itemize}
\raggedright
  \item \textquotedblleft{}Are you travelling tomorrow?\textquotedblright{} … \textquotedblleft{}It depends on the weather\textquotedblright{} — \emph{hal sa-tusāfiru ghadan?} … \emph{yaʿtamidu ʿalā aṭ-ṭaqs}
  \item \textquotedblleft{}It's up to you\textquotedblright{} — \emph{yaʿtamidu ʿalayka}
  \item \textquotedblleft{}According to the timetable, the train leaves in ten minutes\textquotedblright{}
  \item \textquotedblleft{}We'll travel by train or by bus, depending on the weather\textquotedblright{} — \emph{sa-nusāfiru bi-l-qiṭār aw bi-l-ḥāfila, ḥasaba aṭ-ṭaqs}
\end{itemize}

\begin{tcolorbox}[breakable,colback=teal!4,colframe=teal!35!black,title={Before you move on}]
\textquotedblleft{}It depends on the weather\textquotedblright{}? (\emph{\textbf{Yaʿtamidu ʿalā aṭ-ṭaqs.}}) \textquotedblleft{}According to the timetable\textquotedblright{}? (\emph{\textbf{Ḥasaba al-jadwal.}}) What does \emph{yaʿtamidu} leave open that \emph{idhā} closes? (\textbf{Which way it goes.})
\end{tcolorbox}
\section[\texorpdfstring{\ar{لو} \ar{كان} \ar{عندي} \ar{وقت} … (law kāna ʿindī …)}{لو كان عندي وقت … (law kāna ʿindī …)}]{\ar{لو} (law) — if it were so}

\label{lesson:AR-C292-law}



\begin{quote}
Say \textquotedblleft{}it depends on the weather\textquotedblright{}, then \textquotedblleft{}according to the timetable\textquotedblright{}. (\emph{Yaʿtamidu ʿalā aṭ-ṭaqs}, \emph{ḥasaba al-jadwal}.)

\emph{Idhā} is for what may well happen. For what is not so, and you only imagine, Arabic has a third \textquotedblleft{}if\textquotedblright{}.
\end{quote}

\begin{grammarlens}[title={Grammar Lens: law, and la- on the result}]
\begin{quote}
\textbf{\ar{لو} \ar{كان} \ar{عندي} \ar{وقت} \ar{لسافرت}.} — \emph{law kāna ʿindī waqt la-sāfartu} — If I had time, I would travel.
\end{quote}

\textbf{\ar{لو}} — \emph{law} — takes the past, and so does the result. The result usually carries \textbf{\ar{لـ}} — \emph{la-} — written onto it: \emph{la-sāfartu}, I would have travelled, I would travel. The meaning is \textquotedblleft{}but I don't have time\textquotedblright{}. That is the difference from \emph{idhā}: \emph{idhā kāna ʿindī waqt, fa-sa-usāfiru} leaves it open; \emph{law} says it is not so.

You have said \emph{law} since chapter six: \emph{law samaḥt}, \textquotedblleft{}if you would permit\textquotedblright{}, the polite \textquotedblleft{}please\textquotedblright{} that does not take the yes for granted.
\end{grammarlens}

\begin{grammarlens}[title={Grammar Lens: if I were you, and if only}]
\begin{quote}
\textbf{\ar{لو} \ar{كنت} \ar{مكانك} \ar{لذهبت} \ar{إلى} \ar{الطبيب}.} — \emph{law kuntu makānaka la-dhahabtu ilā aṭ-ṭabīb} — If I were you, I'd go to the doctor.
\end{quote}

Word for word, \textquotedblleft{}if I were in your place\textquotedblright{}. \textbf{\ar{ليت}} — \emph{layta}, \textquotedblleft{}if only\textquotedblright{} — wishes without a result: \emph{laytanī kuntu ghaniyyan}, if only I were rich; \emph{laytanī kuntu hunāka}, if only I had been there.
\end{grammarlens}

\subsection*{Your turn}
Say these aloud:

\begin{itemize}
\raggedright
  \item \textquotedblleft{}If I had time, I would travel\textquotedblright{} — \emph{law kāna ʿindī waqt la-sāfartu}
  \item the open one, with \emph{idhā} — \emph{idhā kāna ʿindī waqt, fa-sa-usāfiru}
  \item \textquotedblleft{}If I were you, I'd go to the doctor\textquotedblright{} — \emph{law kuntu makānaka la-dhahabtu ilā aṭ-ṭabīb}
  \item \textquotedblleft{}If only I were rich\textquotedblright{} — \emph{laytanī kuntu ghaniyyan}
  \item a wish of your own with \emph{law} and \emph{la-}, then one with \emph{atamannā an}
\end{itemize}

\begin{tcolorbox}[breakable,colback=teal!4,colframe=teal!35!black,title={Before you move on}]
\textquotedblleft{}If I had time, I would travel\textquotedblright{}? (\emph{\textbf{Law kāna ʿindī waqt la-sāfartu.}}) What goes on the result? (\emph{\textbf{La-.}}) Which \textquotedblleft{}if\textquotedblright{} says \textquotedblleft{}but it is not so\textquotedblright{}? (\emph{\textbf{Law.}}) \textquotedblleft{}If only\textquotedblright{}? (\emph{\textbf{Layta.}})
\end{tcolorbox}
\section[(review)]{Review — would, if only, and it depends}

\label{lesson:AR-R292-would-and-depends}



\begin{quote}
Say \textquotedblleft{}if it rains, we'll watch a film at home\textquotedblright{}. (\emph{Idhā amṭarat, fa-sa-nushāhidu fīlman fī al-manzil}.) Then \textquotedblleft{}if you arrive early\textquotedblright{}, with the other \textquotedblleft{}if\textquotedblright{}. (\emph{In waṣalta mubakkiran}.)
\end{quote}

\begin{grammarlens}[title={What you've built}]
Three lessons, and three distances from the facts. \emph{Idhā}: it may well happen, \emph{idhā kāna ʿindī waqt, fa-sa-usāfiru}. \emph{Yaʿtamidu ʿalā}: it is still open, \emph{yaʿtamidu ʿalā al-waqt}. \emph{Law}: it is not so, \emph{law kāna ʿindī waqt la-sāfartu}; and \emph{layta}, the wish with no result.

All three, on one question:

\begin{quote}
— \emph{hal sa-tusāfiru?} — \emph{yaʿtamidu ʿalā aṭ-ṭaqs. law kāna ʿindī waqt la-sāfartu al-usbūʿa al-qādim.}
\end{quote}
\end{grammarlens}

\subsection*{Your turn}
Say these aloud:

\begin{itemize}
\raggedright
  \item \textquotedblleft{}It depends on the person in charge\textquotedblright{} — \emph{yaʿtamidu ʿalā al-masʾūl}
  \item \textquotedblleft{}If I were you, I would ask the way\textquotedblright{} — \emph{law kuntu makānaka la-saʾaltu: kayfa aṣilu ilā al-maḥaṭṭa?}
  \item \textquotedblleft{}You ought to ask the client\textquotedblright{} — \emph{yanbaghī an tasʾala al-ʿamīl}
  \item \textquotedblleft{}Take the first street on the right\textquotedblright{} — \emph{khudh ash-shāriʿa al-awwala ʿalā al-yamīn}
  \item \textquotedblleft{}If only the sentence were shorter\textquotedblright{} — \emph{layta al-jumlata kānat aqṣar}
  \item \textquotedblleft{}If it were a new experience, I would try it\textquotedblright{} — \emph{law kānat tajribatan jadīdatan la-jarrabtuhā}
  \item one thing you would do \emph{law kāna ʿindī waqt}, and one you will do \emph{idhā} something happens
\end{itemize}

\begin{tcolorbox}[breakable,colback=teal!4,colframe=teal!35!black,title={Before you move on}]
Which \textquotedblleft{}if\textquotedblright{} says it is not so? (\emph{\textbf{Law.}}) What goes on its result? (\emph{\textbf{La-.}}) \textquotedblleft{}It depends on the time\textquotedblright{}? (\emph{\textbf{Yaʿtamidu ʿalā al-waqt.}}) \textquotedblleft{}If only\textquotedblright{}? (\emph{\textbf{Layta.}})
\end{tcolorbox}
\section[(six lines)]{Six lines about a trip}

\label{lesson:AR-W292-sitta-sutur}



\begin{quote}
Chapter forty-five asked for four lines that held together. This one asks for six, and they tell a story.

Pick one reader: a friend, by name.
\end{quote}

\subsection*{What to know first}
A told story is not a list of events. It has a \textbf{scene}, in \emph{kāna} and \emph{kāna} with the present. It has \textbf{events} in the plain past, put in order with \emph{awwalan}, \emph{thumma}, \emph{ʿindamā} and \emph{fī an-nihāya}. It has at least one \textbf{reason}, with \emph{li-anna}, \emph{bi-sabab} or \emph{fa-}, and a \textbf{feeling}, with \emph{kuntu} and \emph{-an}, or \emph{shaʿartu bi-}.

A good one ends by looking forward or wishing: \emph{atamannā an …}, \emph{idhā … fa-sa- …}, or \emph{law … la- …}.

At B1 a connected text runs longer than six lines. Six lines is the first step towards it, not the whole of it.

\subsection*{Writing — connected composition}
On paper, right to left, write \textbf{six lines} to your reader about a trip, real or invented, one line for each of these:

\begin{enumerate}
\raggedright
  \item the scene: when it was, with \emph{qabla} or \emph{al-māḍī}, and the weather, with \emph{kāna};
  \item where you went and who was with you, with \emph{maʿa}, and the ticket: \emph{tadhkirat dhahāb wa-iyāb} or \emph{dhahāb faqaṭ}, and from which platform;
  \item the first thing that happened, with \emph{awwalan} or \emph{ʿindamā};
  \item the next thing, with \emph{thumma} or \emph{baʿda an}, and a reason, with \emph{li-anna}, \emph{bi-sabab}, \emph{as-sababu huwa anna} or \emph{fa-}, or what you had to do, with \emph{kāna yajibu an};
  \item how it ended and how you felt, with \emph{fī an-nihāya} and \emph{kuntu … -an} or \emph{shaʿartu bi-};
  \item one line forward: \emph{atamannā an}, a real \emph{idhā}, or \emph{law … la-}.
\end{enumerate}

Use only words you already write. The whole exercise is choosing the right past for each line, and joining them.

\begin{tcolorbox}[breakable,colback=teal!4,colframe=teal!35!black,title={Before you move on}]
Now read your lines as you read the trip to the sea in chapter two hundred and eighty-four. Mark every verb \textbf{S} for scene or \textbf{E} for event. Line one should be mostly \textbf{S}; lines three and four mostly \textbf{E}. A line about the weather marked \textbf{E} is the commonest slip: change \emph{amṭarat} back to \emph{kāna aṭ-ṭaqs mumṭiran} if you meant the scene.

Then three spelling checks. Every \emph{wa-} and \emph{fa-} is joined to the next word, with no space: \textbf{\ar{وكان}}, \textbf{\ar{فذهبنا}}. Every \emph{-an} after \emph{kāna} on a word without \textbf{\ar{ة}} has its \textbf{\ar{ا}}: \textbf{\ar{سعيدا}}. And each hamza sits on its seat: \textbf{\ar{سؤال}} with \textbf{\ar{ؤ}}, \textbf{\ar{طائرة}} with \textbf{\ar{ئ}}.

Last, read line six alone. If it has \emph{idhā}, is the verb after it in the past? If it has \emph{law}, does the result carry \emph{la-}?
\end{tcolorbox}
